\documentclass[10pt,a4paper]{amsart}
\usepackage{amsmath,amssymb,amsthm,graphicx}
\usepackage{verbatim}
\newenvironment{reference}{\comment}{\endcomment}
%\newenvironment{reference}{}{}
\newenvironment{slogan}{\comment}{\endcomment}
\newenvironment{history}{\comment}{\endcomment}

% For commutative diagrams we use Xy-pic
\usepackage[all]{xy}

% We use 2cell for 2-commutative diagrams.
\xyoption{2cell}
% Theorem environments.
%
\theoremstyle{plain}
\newtheorem{theorem}[subsection]{Theorem}
\newtheorem{proposition}[subsection]{Proposition}
\newtheorem{lemma}[subsection]{Lemma}

\theoremstyle{definition}
\newtheorem{definition}[subsection]{Definition}
\newtheorem{example}[subsection]{Example}
\newtheorem{exercise}[subsection]{Exercise}
\newtheorem{situation}[subsection]{Situation}

\theoremstyle{remark}
\newtheorem{remark}[subsection]{Remark}
\newtheorem{remarks}[subsection]{Remarks}

\numberwithin{equation}{subsection}

% Macros
%
\def\lim{\mathop{\mathrm{lim}}\nolimits}
\def\colim{\mathop{\mathrm{colim}}\nolimits}
\def\Spec{\mathop{\mathrm{Spec}}}
\def\Hom{\mathop{\mathrm{Hom}}\nolimits}
\def\Ext{\mathop{\mathrm{Ext}}\nolimits}
\def\SheafHom{\mathop{\mathcal{H}\!\mathit{om}}\nolimits}
\def\SheafExt{\mathop{\mathcal{E}\!\mathit{xt}}\nolimits}
\def\Sch{\mathit{Sch}}
\def\Mor{\mathop{\mathrm{Mor}}\nolimits}
\def\Ob{\mathop{\mathrm{Ob}}\nolimits}
\def\Sh{\mathop{\mathit{Sh}}\nolimits}
\def\NL{\mathop{N\!L}\nolimits}
\def\CH{\mathop{\mathrm{CH}}\nolimits}
\def\proetale{{pro\text{-}\acute{e}tale}}
\def\etale{{\acute{e}tale}}
\def\QCoh{\mathit{QCoh}}
\def\Ker{\mathop{\mathrm{Ker}}}
\def\Im{\mathop{\mathrm{Im}}}
\def\Coker{\mathop{\mathrm{Coker}}}
\def\Coim{\mathop{\mathrm{Coim}}}

% Boxtimes
%
\DeclareMathSymbol{\boxtimes}{\mathbin}{AMSa}{"02}

%
% Macros for moduli stacks/spaces
%
\def\QCohstack{\mathcal{QC}\!\mathit{oh}}
\def\Cohstack{\mathcal{C}\!\mathit{oh}}
\def\Spacesstack{\mathcal{S}\!\mathit{paces}}
\def\Quotfunctor{\mathrm{Quot}}
\def\Hilbfunctor{\mathrm{Hilb}}
\def\Curvesstack{\mathcal{C}\!\mathit{urves}}
\def\Polarizedstack{\mathcal{P}\!\mathit{olarized}}
\def\Complexesstack{\mathcal{C}\!\mathit{omplexes}}
% \Pic is the operator that assigns to X its picard group, usage \Pic(X)
% \Picardstack_{X/B} denotes the Picard stack of X over B
% \Picardfunctor_{X/B} denotes the Picard functor of X over B
\def\Pic{\mathop{\mathrm{Pic}}\nolimits}
\def\Picardstack{\mathcal{P}\!\mathit{ic}}
\def\Picardfunctor{\mathrm{Pic}}
\def\Deformationcategory{\mathcal{D}\!\mathit{ef}}
\usepackage{cleveref}
\crefname{originalreference}{original reference}{original references}
\Crefname{originalreference}{Original reference}{Original references}
\makeatletter
\newlabel{lemma-normal-closure}{{09DT}{source}{}{sourcepointer-1}{}}
\newlabel{lemma-normal-closure@cref}{{[originalreference][0][]09DT}{[1][0][]source}{}{}{}}
\makeatother
\title{The fundamental theorem of algebra: the complex numbers are algebraically closed}
\author{The Stacks Project Authors}
\date{stacks-project:fields.tex:local-snapshot-20261012; GNU FDL 1.2 or later}
\begin{document}
\maketitle
\noindent\textit{Editorial note.} All delivered mathematical text is the author's own, copied from the pinned Stacks Project chapter file fields.tex; no statement or proof step was rewritten, reordered or removed.\par\medskip
\noindent\textit{Editorial note.} The counted proof is the author's complete argument: The proof uses the intermediate value theorem, i.e. an analytic input, as the author notes in the section opening.\par\medskip
\noindent\textit{Editorial note.} The label lemma-C-algebraically-closed names the counted statement, and the proof environment immediately following it in the source is delivered verbatim as its proof.\par\medskip
\noindent\textit{Editorial note.} The counted statement is carried by the author's own lemma environment in the chapter source.\par\medskip
\noindent\textit{Editorial note.} Revision r002: regenerated from the same spec after the r001 assembly failed at compile time (the r001 context span ran through line 3422 and overlapped the statement span 3418-3421, so the lemma environment and its label were injected twice and the engine reported multiply-defined labels; the context now ends at line 3416, before the statement environment begins; the stale r001 job directory is retained as evidence).\par\medskip
\noindent\textit{Allowed standard prerequisites (editorial).}\begin{itemize}
\item The intermediate value theorem from calculus: every odd degree real polynomial has a real root, by sign change of the dominating leading term
\item Galois theory for finite extensions: for a finite Galois extension with group G and a subgroup P, the fixed field K\textasciicircum{}P has degree over the base equal to the index of P in G
\item The 2-sylow theorem: every finite group has a 2-sylow subgroup, and the index of a 2-sylow subgroup is odd
\item Up to isomorphism the only degree two extension of the reals is the complex numbers, and every complex number has a square root, so the complex numbers admit no quadratic extension
\item Fields, Lemma lemma-normal-closure: every finite extension can be embedded in a finite normal extension over the base
\end{itemize}\medskip
\ifdefined\hypertarget\hypertarget{sourcepointer-1}{}\fi
\noindent\textit{Source reference.} Original reference 09DT: sources/stacks-fields/pinned-original-source.tex, source lines 2303--2303. Prerequisite pointer invoked at 3469; it names a standard fact of the same chapter that the counted argument uses, and every step of the counted proof built on it is retained verbatim in the delivered body.\par\medskip
\section{The complex numbers}
\label{section-complex-numbers}

\noindent
The fundamental theorem of algebra states that the field of complex numbers
is an algebraically closed field. In this section we discuss this
briefly.

\medskip\noindent
The first remark we'd like to make is that you need to use a little
bit of input from calculus in order to prove this. We will use the
intuitively clear fact that every odd degree polynomial over
the reals has a real root. Namely, let
$P(x) = a_{2k + 1} x^{2k + 1} + \ldots + a_0 \in \mathbf{R}[x]$
for some $k \geq 0$ and $a_{2k + 1} \not = 0$.
We may and do assume $a_{2k + 1} > 0$. Then for $x \in \mathbf{R}$
very large (positive) we see that $P(x) > 0$ as the term
$a_{2k + 1} x^{2k + 1}$ dominates all the other terms. Similarly,
if $x \ll 0$, then $P(x) < 0$ by the same reason (and this is where
we use that the degree is odd). Hence by the intermediate value
theorem there is an $x \in \mathbf{R}$ with $P(x) = 0$.

\medskip\noindent
A conclusion we can draw from the above is that $\mathbf{R}$ has
no nontrivial odd degree field extensions, as elements of such extensions
would have odd degree minimal polynomials.

\medskip\noindent
Next, let $K/\mathbf{R}$ be a finite Galois extension with Galois group $G$.
Let $P \subset G$ be a $2$-sylow subgroup. Then $K^P/\mathbf{R}$ is an
odd degree extension, hence by the above $K^P = \mathbf{R}$,
which in turn implies
$G = P$. (All of these arguments rely on Galois theory of course.)
Thus $G$ is a $2$-group. If $G$ is nontrivial, then we see that
$\mathbf{C} \subset K$ as $\mathbf{C}$ is (up to isomorphism) the only
degree $2$ extension of $\mathbf{R}$. If $G$ has more than $2$ elements
we would obtain a quadratic extension of $\mathbf{C}$.
This is absurd as every complex number has a square root.

\medskip\noindent
The conclusion: $\mathbf{C}$ is algebraically closed. Namely, if not
then we'd get a nontrivial finite extension $K/\mathbf{C}$
which we could assume normal (hence Galois) over $\mathbf{R}$ by
Lemma \ref{lemma-normal-closure}. But we've seen above that then
$K = \mathbf{C}$.

\begin{lemma}[Fundamental theorem of algebra]
\label{lemma-C-algebraically-closed}
The field $\mathbf{C}$ is algebraically closed.
\end{lemma}

\begin{proof}
See discussion above.
\end{proof}

\end{document}
